\scanpage{63}
% ll. 1–4
Pour chacun des \inserted{quatre autres} cas $K=-1,0,\alpha,\alpha'$, ceci fait deux autres possibilités, ce qui fait a priori encore $8\cdot15=120$ possibilités --- mais il faut en exclure un bon paquet \ldots !

% ll. 5–9
Principes et critères d'élimination

a) la condition
\[
0<\zeta<1
\]
exprimant que la forme \deleted{\unsure{sesq.}} hermitienne inv. est définie

% ll. 10–16
b) la condition que \inserted{(le $=$)} $\zeta$ satisfasse des équations analogues quand on change $\theta_0$ en $\deleted{\unsure{k}}i\theta_0$ ($i\not\equiv0\ (\nu_0)$) et $\theta_1$ en $j\theta_1$ ($j\not\equiv0\ (\nu_1)$), quitte à changer $K_{\unsure{*}}$ dans $\{\unsure{2|-2},-1,0,\alpha,\alpha'\}$ --- i.e. que l'ens. des valeurs prises par
\[
16\sin^2\frac{\theta}{2}\sin^2\frac{\theta'}{2}
\left(\zeta^2-\frac{\cos\frac{\theta'-\theta}{2}}{\sin\frac{\deleted{\unsure{\theta_0}}\theta}{2}\sin\frac{\deleted{\unsure{\theta_1}}\theta'}{2}}\zeta+
\frac18\frac{\cos(\theta'-\theta)}{\sin^2\frac{\theta}{2}\sin^2\frac{\theta'}{2}}\right)
\]
pour $\theta=i\theta_0,\ \theta'=j\theta_1$ $(i\not\equiv0\ (\nu_0),\ \theta'\not\equiv0\ (\nu_1))$ est contenu dans $\{-2,-1,0,\alpha,\alpha'\}$ \ldots

\begin{center}\tikz\draw (0,0)--(1.4,0)--(.55,-.04).. controls (1.5,-.08) and (.9,-.15)..(1.09,-.35)--(1.16,-.47);\end{center}
% ll. 17–25
L'analyse géométrique via ``\unsure{syst. de Schwartz|syst. de Schwarz}'' montre que la \uline{situation projective} (dans $C=\check{\mathbb P}(V)$) est déterminée mod iso, par la connaissance  des $K_0,K_1$ et $K$ \deleted{\unsure{à}} (que fixent les \deleted{\unsure{trois} $K_0,K_1$} classes de conjugaison de $\tau_0^{C},\tau_1^C$ et $(\tau_0^C\tau_1^C)^{-1}$ (ou $\tau_0^C\tau_1^C$) dans \deleted{D'autre part, pour} $G=\text{\uline{Aut}}(X)$, ou ce qui revient au $=$, dans le groupe \unsure{pythagoricien} $H$ qu'ils engendrent (mod $\operatorname{Aut}(H)$ dans le cas tétraédral, pas mod $\operatorname{Int}(H)$))
% archive sheet number: 62
